%!TEX root = ../main.tex

\section*{Introduction}

Microbial nutrient utilization phenotypes are key determinants of ecological niche and influence organism distribution, interactions, and functional roles within complex ecosystems \cite{Falkowski2008-bx}.
A mechanistic understanding of these phenotypes is crucial for advancing microbial ecology, providing insights into community assembly, nutrient cycling, and biotechnological applications such as probiotic design or bioremediation \cite{Lu2014-zz}.
Traditional laboratory-based cultivation methods have characterized most existing microbial phenotype data, but these methods cannot characterize the metabolic capabilities of unculturable taxa \cite{Lloyd2018-ar}, introducing significant biases and restricting our understanding of microbial diversity.
Even for culturable organisms, comprehensive phenotypic characterization across environmental conditions remains resource-intensive \cite{Barberan2017-xi}.

Rapid advances in genomics offer a promising path to overcome these limitations.
The increasing availability of genomic data, including metagenome-assembled genomes (MAGs), presents an opportunity to predict phenotypes directly from genomic information \cite{Li2023-up}, bypassing extensive laboratory experiments.
While phylogeny has commonly been used as a proxy for phenotype prediction, it is insufficient for characterizing novel organisms with limited phylogenetic context or for traits that are not phylogenetically conserved \cite{Bradley2018-rs}.
Genome-based predictions offer a powerful alternative, enabling characterization of both cultured and uncultured taxa while potentially uncovering the genetic mechanisms driving phenotypic variation.
These predictions can leverage various genomic features, from general characteristics such as GC content \cite{Hu2022-fr} and genome size to specific information derived from gene annotations and inferred metabolic pathways \cite{Long2021-yz}.

Machine learning models have emerged as valuable tools for inferring phenotypes from complex genomic data, capable of handling the non-linear relationships and high dimensionality often encountered in biological systems \cite{Greener2022-dy}.
However, the performance and generalizability of these models are constrained by data limitations and evaluation challenges.
Limited and taxonomically biased training data can lead to overfitting and poor applicability to diverse microbial communities \cite{Ramoneda2024-yb,Albright2023-ze}.
Furthermore, phylogenetic correlations among genomic features confound the identification of mechanistically relevant predictors, making it difficult to distinguish whether models learn true biological relationships or spurious patterns tied to evolutionary history \cite{Li2023-up}.
Previous work has demonstrated these generalization failures and concluded that massive datasets are needed to overcome phylogenetic confounding.
However, an alternative hypothesis---that data quality rather than quantity drives generalization---has received less attention.

We tested this hypothesis using microbial carbon utilization as a model system.
Carbon utilization represents an ideal test case due to its ecological and technological importance, and because these phenotypes involve highly conserved, well-characterized mechanisms---creating conditions where generalizability should be high and interpretability can be rigorously assessed.
We compiled a comprehensive dataset of 819 microbial genomes and their binary growth outcomes across 242 carbon sources from four independent experimental collections.
We developed phenotype-specific classifiers using this dataset and employed various feature representations and evaluation strategies within a standardized framework.

Our analyses revealed that models trained on aggregated data failed to generalize across datasets and phylogenetic contexts, learning dataset-specific correlations rather than mechanistic features.
By stratifying samples according to agreement between mechanistic predictions and experimental outcomes, we found that models trained on high-quality samples achieved robust generalization and consistently identified established pathway genes as top features.
Critically, these models also predicted a subset of samples that the mechanistic predictor misclassified, demonstrating that machine learning can capture relationships beyond curated pathway rules.
We evaluated multiple improvement strategies, including confidence-based filtering for situations where mechanistic predictors are unavailable, and established that modest sample sizes (50--100 high-quality samples per phenotype) are sufficient for robust generalization.
Importantly, we found that full-dataset training remains appropriate when test conditions resemble training conditions, as it maximizes performance in same-environment applications.
Our findings support the hypothesis that data quality, not algorithmic complexity, is the primary constraint on generalizable phenotype prediction, and provide a framework for prioritizing data curation in biological machine learning applications.
